\question[4]
Die Dualdarstellung einer Zahl $n$ erhält man, indem man $n$ wiederholt durch 2 teilt und sich die Reste merkt. Der zuerst
berechnete Rest ist die letzte Ziffer. Die Funktion \texttt{dualzahl(n)} bekommt eine ganze Zahl $n \geq 0$ und gibt
die Dualdarstellung als String zurück. Benutze für die Reste eine Python-Liste als Keller.
Die Funktion \texttt{bin} darf nicht benutzt werden.
\begin{lstlisting}
print(dualzahl(6), dualzahl(37), dualzahl(0))

110 100101 0
\end{lstlisting}
\begin{solutionbox}{5.5cm}
\begin{lstlisting}
def dualzahl(n):
    if n == 0:
        return '0'
    k = []
    while n > 0:
        k.append(n % 2)
        n = n // 2
    s = ''
    while k:
        s += str(k.pop())
    return s
\end{lstlisting}
\end{solutionbox}
